\chapter{First Fractions}\label{ch:g4:fractions}

Half an apple, a quarter of an hour, three quarters of the class:
fractions name the pieces of things. This chapter introduces the
writing $\frac{a}{b}$, reads fractions on pictures and number lines,
and compares the simplest ones. (Fractions grow up in
\cref{ch:g6:fractions}.)

\section{Naming the pieces}

\begin{definition}[Fraction]\label{def:g4:fractions:def}
Cut a whole into equal parts and take some:
\[
\frac{3}{4}
\quad
\begin{array}{l}
3 \text{: how many parts are taken (the \emph{numerator});}\\
4 \text{: how many equal parts the whole is cut into (the
\emph{denominator}).}
\end{array}
\]
$\frac12$ is a \emph{half}, $\frac13$ a \emph{third}, $\frac14$ a
\emph{quarter}, $\frac{1}{10}$ a \emph{tenth}.
\end{definition}

\begin{omfigure}
\begin{tikzpicture}[scale=0.78]
  % half circle
  \draw[thick] (0,0) circle (1.0);
  \fill[omDef!40] (0,0) -- (1,0) arc (0:180:1) -- cycle;
  \draw[thick] (-1,0) -- (1,0);
  \node[below, font=\small] at (0,-1.2) {$\frac12$};
  % quarters
  \begin{scope}[xshift=3.6cm]
  \draw[thick] (0,0) circle (1.0);
  \fill[omThm!40] (0,0) -- (1,0) arc (0:90:1) -- cycle;
  \fill[omThm!40] (0,0) -- (0,1) arc (90:180:1) -- cycle;
  \fill[omThm!40] (0,0) -- (-1,0) arc (180:270:1) -- cycle;
  \draw[thick] (-1,0) -- (1,0);
  \draw[thick] (0,-1) -- (0,1);
  \node[below, font=\small] at (0,-1.2) {$\frac34$};
  \end{scope}
  % thirds bar
  \begin{scope}[xshift=6.6cm, yshift=-0.35cm]
  \foreach \i in {0,1,2} \draw[omExo] (\i*1.3,0) rectangle
    (\i*1.3+1.3,0.7);
  \fill[omMeth!40] (0.04,0.04) rectangle (1.26,0.66);
  \draw[thick] (0,0) rectangle (3.9,0.7);
  \node[below, font=\small] at (1.95,-0.15) {$\frac13$};
  \end{scope}
  % tenths bar
  \begin{scope}[xshift=11.6cm, yshift=-0.35cm]
  \foreach \i in {0,...,9} \draw[omExo] (\i*0.42,0) rectangle
    (\i*0.42+0.42,0.7);
  \foreach \i in {0,...,6} \fill[omProp!40] (\i*0.42+0.03,0.04)
    rectangle (\i*0.42+0.39,0.66);
  \draw[thick] (0,0) rectangle (4.2,0.7);
  \node[below, font=\small] at (2.1,-0.15) {$\frac{7}{10}$};
  \end{scope}
\end{tikzpicture}

{\small Reading fractions on pictures. The parts must be
\emph{equal}: three unequal pieces do not make thirds!}
\end{omfigure}

\begin{example}[Fractions of a collection]\label{ex:g4:fractions:collection}
$\frac13$ of $12$ marbles: share the $12$ marbles into $3$ equal
groups of $4$; one group is $\frac13$, so $\frac13$ of $12$ is $4$
--- and $\frac23$ of $12$ is two groups, $8$.
\end{example}

\section{Fractions on the number line}

\begin{method}[Placing $\frac{a}{b}$]\label{met:g4:fractions:line}
Cut each unit of the line into $b$ equal steps; from $0$, walk $a$
steps. If $a$ is bigger than $b$, the walk passes $1$: fractions can
be bigger than one whole!
\end{method}

\begin{omfigure}
\begin{tikzpicture}[scale=1.6]
  \draw[-Stealth] (-0.2,0) -- (6.9,0);
  \foreach \i in {0,4,8,12,16,20,24}
    \draw ({\i*0.25},-0.09) -- ({\i*0.25},0.09);
  \foreach \i in {0,...,24} \draw ({\i*0.25},-0.05) -- ({\i*0.25},0.05);
  \foreach \x/\l in {0/0, 1/1, 2/2, 3/3, 4/4, 5/5, 6/6}
    \node[below, font=\small] at (\x,-0.1) {$\l$};
  \fill[omDef] (0.5,0) circle (1.6pt)
    node[above, font=\small] {$\frac24$};
  \fill[omThm] (1.25,0) circle (1.6pt)
    node[above, font=\small] {$\frac54$};
  \fill[omMeth] (2.75,0) circle (1.6pt)
    node[above, font=\small] {$\frac{11}{4}$};
\end{tikzpicture}

{\small The line cut in quarters: $\frac24$ lands on the same point
as $\frac12$; $\frac54$ is one quarter past $1$; $\frac{11}{4}$ is
$2$ and $\frac34$.}
\end{omfigure}

\begin{example}[Whole numbers hiding in fractions]\label{ex:g4:fractions:whole}
$\frac44 = 1$ (four quarters make the whole); $\frac82 = 4$;
$\frac{12}{3} = 4$. And $\frac74$ is $1 + \frac34$: one whole and
three quarters more.
\end{example}

\section{Comparing and adding simple fractions}

\begin{proposition}[Same denominator]\label{prop:g4:fractions:compare}
With the same denominator, the pieces have the same size --- so just
count them:
\[
\frac{2}{8} < \frac{5}{8},
\qquad
\frac{3}{8} + \frac{4}{8} = \frac{7}{8}.
\]
Comparing to one whole is also easy: $\frac{5}{8} < 1 < \frac{9}{8}$
(fewer, then more, than $8$ eighths).
\end{proposition}

\begin{proof}[Why, on a picture]
Shading $3$ eighths and then $4$ more eighths of the same bar shades
$7$ eighths in total.
\end{proof}

\begin{example}[Different denominators, same point]\label{ex:g4:fractions:equal}
On the number line, $\frac12$, $\frac24$ and $\frac{5}{10}$ all land
on the same point: they are three names of the same number. Cutting
each half into two makes quarters; into five, tenths. (The general
rule is in \cref{ch:g6:fractions}.)
\end{example}

\begin{example}[Quarters of an hour]\label{ex:g4:fractions:hour}
An hour has $60$ minutes, so a quarter of an hour is
$60 \div 4 = 15$ minutes, and three quarters of an hour is
$45$ minutes. ``Half past two'' is $\frac12$ hour after two
o'clock.
\end{example}

\section{Exercises}

\begin{exercise}[$\star$]\label{exo:g4:fractions:1}
Write the fraction shown: a pizza cut in $6$ with $5$ slices left;
a chocolate bar of $8$ squares with $3$ eaten (fraction eaten); a
flag divided in $3$ equal vertical bands, $1$ colored.
\end{exercise}

\begin{exercise}[$\star$]\label{exo:g4:fractions:2}
Draw a bar cut into $5$ equal parts and shade $\frac35$. Then shade
$\frac{2}{5}$ of another equal bar. Which shading is bigger?
\end{exercise}

\begin{exercise}[$\star$]\label{exo:g4:fractions:3}
Compute: $\frac12$ of $18$; \quad $\frac14$ of $20$; \quad
$\frac13$ of $21$; \quad $\frac34$ of $20$ (three groups of a
quarter).
\end{exercise}

\begin{exercise}[$\star$]\label{exo:g4:fractions:4}
Place on a number line cut in thirds: $\frac13$; \ $\frac33$; \
$\frac53$; \ $2$. Which fraction lands exactly on $2$?
\end{exercise}

\begin{exercise}[$\star$]\label{exo:g4:fractions:5}
Copy and complete with $<$, $>$ or $=$:
\[
\frac38 \;?\; \frac68, \qquad
\frac55 \;?\; 1, \qquad
\frac74 \;?\; 1, \qquad
\frac12 \;?\; \frac24 .
\]
\end{exercise}

\begin{exercise}[$\star$]\label{exo:g4:fractions:6}
Compute: $\frac26 + \frac36$; \quad $\frac58 - \frac28$; \quad
$\frac14 + \frac24$; \quad $1 - \frac13$ (how many thirds make
$1$?).
\end{exercise}

\begin{exercise}[$\star$]\label{exo:g4:fractions:7}
Write as a whole number plus a fraction smaller than $1$:
$\frac54$; \quad $\frac73$; \quad $\frac{9}{2}$.
\end{exercise}

\begin{exercise}[$\star$]\label{exo:g4:fractions:8}
How many minutes are: half an hour; a quarter of an hour; three
quarters of an hour; $\frac13$ of an hour?
\end{exercise}

\begin{exercise}[$\star$]\label{exo:g4:fractions:9}
In a class of $24$ students, $\frac14$ wear glasses. How many
students wear glasses? How many do not?
\end{exercise}

\begin{exercise}[$\star\star$]\label{exo:g4:fractions:10}
Nina ate $\frac38$ of a pizza and Sam ate $\frac48$ of the same
pizza. What fraction did they eat together? What fraction is left?
Who ate more?
\end{exercise}

\begin{exercise}[$\star\star$]\label{exo:g4:fractions:11}
True or false, with a picture or an example: ``$\frac12$ of a small
pizza is less than $\frac14$ of a very big pizza is possible''.
What must one always know before comparing two fractions of
something?
\end{exercise}
